\chapter{选择与诊断：把分数变成可靠判断}
\label{ch:diagnostics}

一个可复用的建模方案包含预处理、模型、超参数、停止规则与决策规则。
本章接续可信评估，把搜索、曲线、概率与不确定性放回这个完整流程。

\begin{keyidea}[分数越好，越需要问清它怎样得到]
选择次数、数据用途、指标方向与估计误差共同影响结论。
在测试集上反复试出一个好结果，不能替代预先约定的验证流程。
\end{keyidea}

\section{参数、超参数与搜索预算}
\label{sec:diag-search}
参数（Parameter）通常由拟合学习，如回归系数与网络权重；
超参数（Hyperparameter）控制模型或学习流程，如正则强度、树深度、学习率。
超参数可以由验证搜索自动选择，不能简单理解为“只有人工才设置”。

\begin{figure}[htbp]
\centering\begin{tikzpicture}[x=1cm,y=1cm]
\draw[->,Quiet] (0,0)--(4,0) node[right,lab] {$\log\lambda$};
\draw[->,Quiet] (0,0)--(0,3.2) node[above,lab] {超参数2};
\foreach \x in {.5,1.8,3.1}{\foreach \y in {.5,1.5,2.5}{\fill[Teal] (\x,\y) circle(2pt);}}
\node[lab] at (1.8,-.7) {网格：9组，只试3个$\lambda$值};
\begin{scope}[xshift=7cm]
\draw[->,Quiet] (0,0)--(4,0) node[right,lab] {$\log\lambda$};
\draw[->,Quiet] (0,0)--(0,3.2) node[above,lab] {超参数2};
\foreach \x/\y in {.3/.6,.7/2.6,1.0/1.4,1.5/2.1,1.9/.3,2.2/1.7,2.7/2.8,3.1/1.1,3.5/.7}{\fill[Amber] (\x,\y) circle(2pt);}
\node[lab] at (1.8,-.7) {随机：9组，覆盖更多$\lambda$值};
\end{scope}
\end{tikzpicture}

\caption{等预算的独立示意。随机搜索能试更多不同坐标值，但不保证每次更优。}
\par\smallskip
\readfigure{网格搜索枚举指定笛卡尔积；3个正则强度与4个树深度共12组，
五折评估需60次折内拟合，最后全开发集重拟合另计一次。不能按训练分数挑选。}
\end{figure}

随机搜索从预先指定分布抽样；若正值跨多个数量级，可考虑对数尺度，
例如$\log_{10}\lambda\sim U[-4,0]$，而非默认在线性区间均匀取值。
两者给不同尺度的搜索预算，不存在通用最优搜索分布。

\quick{先定指标方向、范围、预算、划分与随机种子。候选失败和实际开销也应记录；搜索只在指定范围内比较。}

\section{嵌套交叉验证：评估选择流程}
\label{sec:diag-nested}
普通交叉验证可用于选择超参数；同一组折上最优候选的分数，已经受到选择影响。
若需要估计整套选择流程的泛化表现，可用嵌套交叉验证（Nested CV）。

\begin{figure}[htbp]
\centering\begin{tikzpicture}[node distance=9mm]
\node[box,text width=24mm] (train) {外层训练集};
\node[box,right=of train,text width=32mm] (inner) {内部交叉验证\\选择完整方案};
\node[box,right=of inner,text width=26mm] (refit) {用外层训练集\\重新拟合};
\node[box,right=of refit,text width=24mm] (score) {外层留出折\\只作评估};
\draw[flow] (train)--(inner);\draw[flow] (inner)--(refit);\draw[flow] (refit)--(score);
\node[lab,below=11mm of inner,text width=70mm,align=center] {每个外层折都重新执行选择；\\外层分数评价的是整套选择流程。};
\end{tikzpicture}

\caption{外层评估、内层选择；预处理始终与模型一起在相应训练部分拟合。}
\par\smallskip
\readfigure{外层留出折不参与内层搜索；在不同外层折上选出的超参数可以不同。
汇总外层结果，评价的是“给定新训练样本后执行这套选择流程”的表现。}
\end{figure}

若外层5折、内层3折、每次比较8组候选，内层拟合共$5\times3\times8=120$次；
每个外层选定后重新拟合一次，共5次，因此外层评估阶段125次。
最后在全开发集上再选参并拟合用于部署，属于另外的计算与模型。

\begin{pitfall}[外层也不能无限反复优化]
看到外层结果后持续改变特征、候选范围或搜索策略，会再次产生适应性选择。
嵌套CV不是泄漏或研究者选择的万能修补。
独立测试集可作为另一种评估设计；按数据、用途与预算选择，不要求所有项目都同时堆叠两者。
\end{pitfall}

时间序列与分组数据的内外层都需保持相应边界。
也不要把五个外层模型中分数最高的一个直接当作“最终最佳模型”。

\section{三种曲线：先读横轴，再谈诊断}
\label{sec:diag-curves}
“学习曲线”有时泛指训练过程曲线，阅读时务必先检查横轴。
这里区分训练样本量、一个超参数与训练步数三类图。

\begin{figure}[htbp]
\centering\begin{tikzpicture}
\begin{axis}[width=71mm,height=61mm,xmin=1,xmax=5,ymin=0,ymax=1,
 xlabel={训练样本量（相对）},ylabel={误差，越低越好},xtick={1,3,5},ytick={0,.5,1},title={学习曲线：固定方案},legend style={at={(.5,-.3)},anchor=north,draw=none,font=\small}]
\addplot[Teal,thick,mark=*] coordinates {(1,.08)(2,.13)(3,.17)(4,.19)(5,.20)};\addlegendentry{训练}
\addplot[Amber,thick,mark=square*] coordinates {(1,.80)(2,.52)(3,.39)(4,.32)(5,.28)};\addlegendentry{验证}
\end{axis}
\begin{axis}[at={(78mm,0)},anchor=south west,width=71mm,height=61mm,xmin=1,xmax=5,ymin=0,ymax=1,
 xlabel={模型复杂度（相对）},xtick={1,3,5},ytick={0,.5,1},title={验证曲线：固定样本量},legend style={at={(.5,-.3)},anchor=north,draw=none,font=\small}]
\addplot[Teal,thick,mark=*] coordinates {(1,.75)(2,.4)(3,.2)(4,.1)(5,.03)};\addlegendentry{训练}
\addplot[Amber,thick,mark=square*] coordinates {(1,.80)(2,.48)(3,.30)(4,.4)(5,.65)};\addlegendentry{验证}
\end{axis}
\end{tikzpicture}

\caption{独立构造的趋势示意，不是实验结果；这类形状并非必然定律。}
\par\smallskip
\readfigure{左图改变训练样本量、固定方案，每个规模重新拟合；右图改变复杂度、固定样本量。
纵轴均为误差，越低越好。若换成准确率等得分，方向应对应改变。}
\end{figure}

\begin{center}
\begin{tabularx}{\textwidth}{P{33mm}Y}
\toprule
横轴 & 可以帮助判断什么\\
\midrule
训练样本量 & 在固定流程下，增加同分布数据是否仍可能改善验证表现。\\
超参数 & 候选范围内的表现与敏感性，不保证所有参数都有简单U形关系。\\
Epoch或更新步数 & 优化过程、训练/验证分离和停止时刻；不是样本量变化。\\
\bottomrule
\end{tabularx}
\end{center}

训练与验证都差：先查标签、指标、实现与优化，再考虑容量不足。
训练好、验证差：检查泄漏、划分、分布变化、正则与数据量。
曲线只提供线索，不能单凭间距唯一诊断原因；验证误差也可能暂时低于训练误差。

\section{概率预测：Brier与对数损失}
\label{sec:diag-probability}
二分类模型给出的$p_i$应对应指定正类，$y_i\in\{0,1\}$。
Brier分数（Brier Score）与二元对数损失都直接评价概率，越小越好：
\begin{align}
 \mathrm{BS}&=\frac1n\sum_i(p_i-y_i)^2,\\
 \mathrm{LogLoss}&=-\frac1n\sum_i[y_i\log p_i+(1-y_i)\log(1-p_i)].
\end{align}
对数用自然底；错误且极度自信的概率会受较大惩罚。
若真实类别被赋予零概率，对数损失趋于无穷；软件裁剪概率需说明约定。

\begin{center}
\begin{tabular}{lrrr}
\toprule
模型（同一8样本） & 正类少组的$p$ & 正类多组的$p$ & Brier\\
\midrule
A & 0.1 & 0.9 & 0.2100\\
B & 0.25 & 0.75 & 0.1875\\
\bottomrule
\end{tabular}
\end{center}
第一组标签$[0,0,0,1]$，第二组$[0,1,1,1]$。
阈值0.5时，两模型准确率均为0.75，AUC也均为0.75，但概率误差不同。
B的对数损失约0.562，A约0.655；排序、硬标签与概率质量是不同问题。

若同一条件下$Y\sim\mathrm{Bernoulli}(q)$，则
\[
 \mathbb E[(p-Y)^2]=(p-q)^2+q(1-q).
\]
期望Brier在$p=q$处最小；并不要求对每个随机结果都给出0或1。
二分类这个定义范围为$[0,1]$；多分类若对全部类别平方差求和，范围及缩放约定不同。

\quick{Brier反映整体概率误差，不能单独证明校准完美；比较时保持样本、类别定义与缩放口径一致。}

\section{校准、排序与阈值分工}
\label{sec:diag-calibration}
校准（Calibration）问：预测概率相近的一组样本，其实际正类频率是否相近？
可靠性图（Reliability Diagram）把组内平均预测概率与组内实际频率作比较。

\begin{figure}[htbp]
\centering\begin{tikzpicture}
\begin{axis}[width=85mm,height=75mm,xmin=0,xmax=1,ymin=0,ymax=1,
 xlabel={组内平均预测概率},ylabel={组内实际正类比例},xtick={0,.25,.5,.75,1},ytick={0,.25,.5,.75,1},
 legend style={at={(1.04,.6)},anchor=west,draw=none,font=\small}]
\addplot[Quiet,dashed,domain=0:1]{x};\addlegendentry{对角参考线}
\addplot[Amber,thick,mark=square*] coordinates {(.1,.25)(.9,.75)};\addlegendentry{较过度自信}
\addplot[Teal,only marks,mark=*,mark size=3pt] coordinates {(.25,.25)(.75,.75)};\addlegendentry{本例组频率}
\end{axis}
\end{tikzpicture}

\caption{沿用前例两个各4样本的组。有限分组图只用于说明，不证明总体已经校准。}
\par\smallskip
\readfigure{A将两组概率推向0与1，比本例观察频率更自信。
按组频率替换为B后，排序与0.5阈值标签不变，Brier下降。实际应用不能用测试标签这样修改预测。}
\end{figure}

Sigmoid/Platt或保序回归等校准器，应使用模型拟合之外的校准数据，
或合适的交叉拟合流程；模型、校准器、阈值都确定后再测独立测试集。
分箱数量和少量样本会影响图形，分布变化后原校准关系也可能失效。

若$p$是适用场景的正类后验概率，正确决策代价为0，误报代价$C_{FP}$、漏报代价$C_{FN}$，
两种错误代价非负且不同时为0，
则预测正类的期望代价为$C_{FP}(1-p)$，负类为$C_{FN}p$。
因此可按
\[
 p\geq\frac{C_{FP}}{C_{FP}+C_{FN}}
\]
判正（相等时的约定另定）。例如$C_{FP}=1,C_{FN}=4$给阈值0.2。
这个结论依赖概率与代价假设，不可对任意未校准得分照搬，也不等于任何F1最优阈值。

\clearpage
\section{Platt校准：给固定得分再拟合一个概率映射}
\label{sec:diag-platt}
分类器的得分可以很好地排序，却不一定对应事件频率。
二分类Platt校准用两个参数将固定得分$s(\vct x)$映到概率：
\[
 p_{\mathrm{cal}}(\vct x)=\sigma(a\,s(\vct x)+b).
\]
这里采用$+as+b$的Sigmoid约定；也有文献写成$1/(1+e^{As+B})$，相应参数符号相反。
输入是事先定义的得分，不要求它原本就是概率。

\begin{figure}[H]
\centering\begin{tikzpicture}[node distance=7mm]
\node[box,text width=36mm] (train) {训练集\;T\\拟合基础模型};
\node[box,text width=48mm,right=of train] (cal) {校准集\;C\\固定模型产生得分\\拟合映射$a,b$};
\node[box,text width=36mm,right=of cal] (test) {测试集\;E\\固定完整流程\\一次独立评价};
\draw[flow] (train)--(cal);\draw[flow] (cal)--(test);
\end{tikzpicture}

\caption{简化的三个数据角色；交叉拟合可提高数据利用率，但也必须保持每条校准预测的拟合边界。}
\end{figure}
\readfigure{先训练基础模型，再冻结它并在校准集上产生得分，利用校准标签拟合$a,b$。
测试只评价已经固定的完整流程。基础模型若重新拟合，得分分布可能变化，原校准映射也需重新检查。}

沿用前面的8样本例子：两组标签分别为$[0,0,0,1]$与$[0,1,1,1]$，
基础概率为0.1与0.9。若以其logit为得分，则$s=-\log9,+\log9$。
取$a=1/2,b=0$后，$p_{\mathrm{cal}}=0.25,0.75$，正好对应两组观察频率。
这个小算例解释映射，不把校准数据上的改善当作独立测试收益。

\begin{center}
\begin{tabularx}{\textwidth}{P{31mm}Y}
\toprule
方法 & 灵活性与限制\\
\midrule
Sigmoid / Platt & 两参数，数据需求相对较低；真实校准关系若偏离此形状，仍可能不合适。\\
保序回归 Isotonic & 拟合单调的分段映射；更灵活，也更需防止小校准集过拟合与并列概率。\\
\bottomrule
\end{tabularx}
\end{center}

\begin{pitfall}[校准不创造新的排序信息]
$a>0$时该映射严格递增，精确数学下保留排序和ROC-AUC；$a<0$会反转，$a=0$会压成常数。
实际接口是否限制方向、裁剪或产生并列需查实现。概率质量改变不等于AUC必须提高。
少数类太少或部署分布变化时，应重新检查不确定性与校准适用性。
\end{pitfall}

二元形式不能不加说明地当作多分类联合概率；逐类映射后还需明确归一化与评价方式。
校准器和阈值也属于选择流程，应一并纳入验证。


\section{MCC：同时使用四种计数}
\label{sec:diag-mcc}
马修斯相关系数（Matthews Correlation Coefficient, MCC）是二元真实标签与预测标签的相关系数。
在四个边际因子均非零时，
\[
 \mathrm{MCC}=\frac{TP\,TN-FP\,FN}
 {\sqrt{(TP+FP)(TP+FN)(TN+FP)(TN+FN)}}.
\]
它在$[-1,1]$内，1表示完全一致，$-1$表示完全反转，0表示在此二元编码下没有线性相关。
MCC不是准确率，也不是保证公平或校准的综合评分。

\begin{center}
\begin{tabular}{lrr}
\toprule
100样本例子 & 准确率 & MCC\\
\midrule
$TP=18,FN=12,FP=7,TN=63$ & 0.81 & 约0.529\\
完全正确，且两类均出现 & 1 & 1\\
完全反转，且两类均出现 & 0 & $-1$\\
全部预测为负类 & 取决于真实负类比例 & 原公式分母为0\\
\bottomrule
\end{tabular}
\end{center}

若真实标签或预测标签恒定，相关系数的普通定义失效。
软件可能使用返回0等扩展约定，应记录这个约定，不能把0自动解释成测得了独立性。
多分类MCC另有推广公式，不能直接将多分类矩阵四个位置代入本节二分类式。

\begin{keyidea}[指标服务于问题]
MCC可补充类别不平衡时的标签评价；仍应看混淆矩阵、各类表现与错误代价。
需要概率质量时看概率指标，需要低误报运行点时检查相应阈值表现。
\end{keyidea}

\section{不确定性：区间也有适用条件}
\label{sec:diag-uncertainty}
有限测试样本得到的准确率只是估计。
在模型已固定、测试结果近似独立且同分布的二元成功/失败设置下，
可用二项比例区间描述抽样不确定性。

\begin{figure}[htbp]
\centering\begin{tikzpicture}
\begin{axis}[width=128mm,height=39mm,xmin=.6,xmax=1,ymin=0,ymax=1,
 xlabel={总体准确率的估计与Wilson近似区间},ytick=\empty,xtick={.6,.7,.8,.9,1},axis y line=none,axis x line=bottom]
\draw[Teal,very thick] (axis cs:.7222,.45)--(axis cs:.8749,.45);
\draw[Teal,thick] (axis cs:.7222,.35)--(axis cs:.7222,.55);
\draw[Teal,thick] (axis cs:.8749,.35)--(axis cs:.8749,.55);
\addplot[only marks,Teal,mark=*,mark size=3pt] coordinates {(.81,.45)};
\node[lab,above] at (axis cs:.81,.55) {81/100，约$[0.722,0.875]$};
\end{axis}
\end{tikzpicture}

\caption{100个独立测试结果中81个正确，采用$z=1.96$的Wilson近似95\%区间。}
\par\smallskip
\readfigure{0.81不是没有误差的总体常数。区间约为$[0.722,0.875]$，
这不包含对未来分布变化的保证，也没有自动计入重新训练或调参的不确定性。}
\end{figure}

令$\hat p=k/n$，Wilson区间的中心和半宽分别为
\[
 c=\frac{\hat p+z^2/(2n)}{1+z^2/n},\qquad
 r=\frac{z\sqrt{\hat p(1-\hat p)/n+z^2/(4n^2)}}{1+z^2/n}.
\]
区间为$[c-r,c+r]$；这是近似方法，不对每种有限样本设置保证精确95\%覆盖。
频率学派的95\%指在相应重复抽样下，该构造方法的长期覆盖率，
不是“固定真参数有95\%概率落在这一次固定区间内”。

\begin{pitfall}[不能把任何波动除以根号样本数就叫置信区间]
同一病人的多条记录不是独立试验；交叉验证各折也不相互独立。
折间标准差不等于均值的标准误；模型比较还应保留同一测试样本上的成对关系。
有组结构时按组设计估计或重采样，明确抽样单位与所覆盖的不确定性。
\end{pitfall}

\begingroup\small
\noindent\textbf{参考：}ISLR第二版第5章；Bergstra与Bengio，\textit{Random Search for Hyper-Parameter Optimization}（2012）；
scikit-learn Model Selection、Calibration与Model Evaluation文档；Wilson（1927）。
全部图与算例独立构造。
\par\endgroup
